\documentclass[10pt,a4paper]{amsart}
\usepackage[margin=19mm]{geometry}
\usepackage{amsmath,amssymb,mathtools}
\usepackage[scr=rsfs,cal=euler]{mathalfa}
\usepackage{xfrac}
\usepackage[unicode,hidelinks,bookmarks=false]{hyperref}
\newcommand{\ie}{i.e.\ }
\newcommand{\cf}{cf.\ }
\newcommand{\resp}{respectively\ }
\newcommand{\Subsection}{\subsection}
\providecommand{\nobreakdash}{\nobreak\hbox{-}}
\setlength{\emergencystretch}{2em}
\multlinegap0pt
\usepackage{amssymb}
\usepackage{csquotes}
\usepackage{tikz}
\newtheorem{lemma}{Lemma}
\newcommand{\R}{\mathbb{R}} %réels
\newcommand{\Tsf}{\mathsf{T}} %Transposée
\newcommand{\lV}{\left\Vert} %crochet angle gauche
\newcommand{\lp}{\left(} %parenthèse gauche
\newcommand{\lv}{\left\vert} %crochet angle gauche
\newcommand{\rV}{\right\Vert} %crochet angle droit
\newcommand{\rp}{\right)} %parenthèse droite
\newcommand{\rv}{\right\vert} %crochet angle droit
\title{Dirac resonances as non-self-adjoint eigenvalues}
\author{Henry DUMANT}
\date{Source: arXiv:2607.26166v1 (CC BY 4.0)}
\begin{document}
\maketitle

\noindent\emph{Source and licence.} Dirac resonances as non-self-adjoint eigenvalues. Version arXiv:2607.26166v1 (CC BY 4.0), distributed by arXiv under the Creative Commons Attribution 4.0 International licence, http://creativecommons.org/licenses/by/4.0/, as recorded in the arXiv licence metadata for this exact version and shown on the abstract page https://arxiv.org/abs/2607.26166v1. Full licence notice: \texttt{licenses/arxiv-2607.26166-CC-BY-4.0.txt}.\par\smallskip

\noindent\textit{Editorial scope.} Counted unit: the article's lemma FChoice (source0.tex lines 815--828). The article's complete original proof of it is delivered with it (source0.tex lines 830--882, one \texttt{proof} environment of 53 lines). No mathematics is written, repaired or re-derived: the statement and the proof are the article's own lines, in the article's own notation, and no retained line is substituted or re-flowed.  No further source-local context is needed: every cross-reference of every retained line resolves inside the two delivered spans.  Nothing was omitted from any retained span, so the package carries no omission record.  No retained line contains a citation, so the wrapper carries no bibliography and no bibliography entry.  No retained line contains a figure environment, an image-inclusion command or drawing code, so no image is delivered and the package declares no asset dependency.  Editorial formatting only, in addition: an AMS \texttt{amsart} preamble that restates the article's own declarations and macros used by the retained lines, namely the environments \texttt{lemma} on separate counters, together with the standard packages those lines need.  The retained environments print in this document as Lemma 1 in order of appearance, whereas the article prints the same items with the numbering of its own full document.  The complete native source of the article as distributed in the arXiv e-print for this exact version is bundled unchanged as \texttt{sources/206242/source0.tex}.
\begin{lemma}\label{FChoice}
    For all compact set $K\subset\R^3$ and $\ell>1$ there exists $F\in C^2(\R^3,\R^3)$ satisfying
    \begin{equation*}
        \left\{
        \begin{array}{ll}
            F=0 \mbox{ on a neighborhood of } K\\

            F(x)=x \mbox{ if } \lv x\rv\gg 1\\
            
            \lV \nabla F\rV_{L^\infty}\leq\ell
        \end{array}
        \right.
    \end{equation*}
\end{lemma}

\begin{proof}
    Let $r_0>0$ such that $K\subset B(0,r_0)$. We set $F(x)=f(\lv x\rv)x$ where $f:\R\rightarrow[0,1]$ is $C^2$ and
    \begin{equation*}
        f(r):=\left\{
        \begin{array}{ll}
            0 &\mbox{ if } r\in (-\infty,r_0)\\
            
            1 &\mbox{ if } r\in(r_1,+\infty)
        \end{array}.
        \right.
    \end{equation*}
    for some $r_1>0$ chosen later. We compute for $x\not=0$
    \begin{equation*}
        \nabla F(x)=f(\lv x\rv)I_3+\frac{f'(\lv x\rv)}{\lv x\rv}xx^\Tsf
    \end{equation*}
    therefore 
    \begin{equation*}
        \lV\nabla F\rV_{L^\infty}\leq 1+\sup_{r\geq0}\lv rf'(r)\rv.
    \end{equation*}
    The problem reduces to the question : can we define $f$ on $[r_0,r_1]$ in such a way that 
    \begin{equation*}
        \sup_{r\geq0}\lv rf'(r)\rv\leq\ell-1\mbox{ ?}
    \end{equation*}
    This constraint indicates a logarithmic behaviour. To preserve smoothness, we introduce a $C^2-$smooth function $g:\R\rightarrow[0,1]$ such that
    \begin{equation*}
         g(r):=\left\{
        \begin{array}{ll}
            0 &\mbox{ if } r\in (-\infty,1/4)\\
            
            1 &\mbox{ if } r\in(3/4,+\infty)
        \end{array},
        \right.
    \end{equation*}
    and for $r\in[r_0,r_1]$ we set
    \begin{equation*}
        f(r):= g\lp\frac{\log r-\log r_0}{\log r_1-\log r_0}\rp.
    \end{equation*}
    That way, we get 
    \begin{equation*}
        f'(r)=\left\{
        \begin{array}{ll}
            0 &\mbox{ if } r\in (-\infty,r_0)\cup(r_1,+\infty)\\
            
            \dfrac{1}{r(\log r_1-\log r_0)} g'\lp\dfrac{\log r-\log r_0}{\log r_1-\log r_0}\rp &\mbox{ if } r\in[r_0,r_1]
        \end{array}
        \right.
    \end{equation*}
    therefore 
    \begin{equation}\label{LogarithmicUpper}
        \sup_{r\geq0}\lv rf'(r)\rv\leq\frac{\lV g'\rV_{L^\infty}}{\log r_1-\log r_0}.
    \end{equation}
    It remains to take $r_1$ large enough so that the r.h.s of \eqref{LogarithmicUpper} is smaller than $\ell-1>0$.
\end{proof}

\end{document}
